% Problem dossier: VI.05.P6
\subsection*{Problem VI.05.P6 --- Three coordinate axes: components, graph and ring}
\label{prob:vi05-p6}
\paragraph{Problem.}
Let
\[
X=V(xy,xz,yz)\subseteq\mathbb A_k^3.
\]
Determine all irreducible components and their prime ideals.  Compute the component graph,
prove that \(X\) is connected, and explain algebraically why \(k[X]\) has zero divisors but
no nontrivial idempotents when \(k\) is algebraically closed.
\ifdefined\FullProblemDossiers
\paragraph{Why this problem matters.}
It combines four dictionaries at once: equations to components, minimal primes to geometry,
component incidence to connectedness, and zero divisors/idempotents to different kinds of
splitting.
\paragraph{Useful ideas before starting.}
A point satisfying \(xy=xz=yz=0\) has at most one nonzero coordinate.  The three coordinate
axes have prime ideals \((y,z)\), \((x,z)\), and \((x,y)\).
\paragraph{Guided hints.}
\textbf{Hint 1.} Prove set-theoretically that every point lies on one coordinate axis.
\textbf{Hint 2.} Check that the three axis ideals are prime and none contains another.
\textbf{Hint 3.} Every pair of axes meets at the origin.
\paragraph{Solution.}
The equations force at most one of \(x,y,z\) to be nonzero, hence
\[
X=V(y,z)\cup V(x,z)\cup V(x,y).
\]
Each coordinate ring is a polynomial ring in one variable, so all three component ideals are
prime.  None of the three axes contains another, hence they are exactly the irreducible
components.  Equivalently,
\[
\sqrt{(xy,xz,yz)}=(y,z)\cap(x,z)\cap(x,y),
\]
and the ideal on the left is already radical.
Every pair of components meets at \((0,0,0)\), so the component graph is a triangle and is
connected.  Thus \(X\) is connected.  The classes of \(x\) and \(y\) are nonzero zero
divisors because \(xy=0\), so the ring is not a domain.  Over algebraically closed \(k\),
connectedness implies that the coordinate ring has no nontrivial idempotents.
\paragraph{What happened mathematically.}
The example makes the hierarchy visible: three irreducible pieces can form one connected
object when the incidence graph ties them together.  Zero divisors see the three branches;
idempotents would require a genuine disjoint separation.
\paragraph{Extensions.}
Generalize to the union of the \(n\) coordinate axes in \(\mathbb A^n\).  Determine the
component graph and the ideal generated by all pairwise products \(x_ix_j\).
\paragraph{Provenance.}
New canonical synthesis problem, related to the legacy unions-of-lines material that is
reserved for tangent-space analysis in VI/32.
\fi
